% paper/sections/10_4_verification_summary.tex
%
% §10.4  検証結果のまとめ
% ═══════════════════════════════════════════════════════════════════════════════
\subsection{検証結果のまとめ}
\label{sec:verify_summary}
% ═══════════════════════════════════════════════════════════════════════════════

表~\ref{tab:verification_summary} に各コンポーネントの検証結果を総括する．
主要な知見を以下に整理する．

\begin{enumerate}
  \item \textbf{空間離散化：}
        CCD 微分は周期境界テストで $\Ord{h^{6.0}}$ を達成し，設計精度を確認した（表~\ref{tab:ccd_periodic}）．
        DCCD フィルタは Nyquist モードを完全に抑制する（$H(\pi)=0$）．
        非一様格子（$\alpha=2$）でも $\Ord{h^6}$ が保全され，
        GCL は計算機精度で成立する（§\ref{sec:gcl_verification}）．
  \item \textbf{DCCD 移流特性：}
        1次元移流比較により，DCCD は不連続プロファイルの TV を CCD の約 $1/3$ に低減し，
        平滑プロファイルでは CCD と同等の精度を維持する（§\ref{sec:dccd_dissipation_bench}）．
  \item \textbf{曲率計算：}
        $\phi$ 直接微分（経路 A）では $\Ord{h^5}$ の高次収束を示すが，
        採用経路 B（$\psi \to \phi$ ロジット逆変換）は界面厚み $\varepsilon = \Ord{h}$ が精度の上界を規定する．
        最近接点 Hermite 補間（§~\ref{sec:extension_pde_disc}）により
        界面越し場延長は $\Ord{h^6}$ を達成するため，界面近傍の律速は
        CSF 正則化 $\delta_\varepsilon$（$\Ord{h^2}$）となり，曲率計算経路の選択は実用上の精度に影響しない．
  \item \textbf{界面処理：}
        CLS 保存型再マッピングは LS 非保存補間に比べ質量誤差を 1--2 桁低減した
        （§\ref{sec:cls_conservation}）．
        上風差分ベースライン（Aslam 2004）は場延長収束テストで $\Ord{h^{1.0}}$ を確認した．
        最近接点 Hermite 補間（§~\ref{sec:extension_pde_disc}）への切り替えにより $\Ord{h^6}$ が達成される．
        Young--Laplace 静止液滴テストでは上風・Hermite 両パイプラインで完全一致し，
        $N = 128$ で相対誤差 $0.2\%$ を達成した（§\ref{sec:verify_field_extension}）．
  \item \textbf{PPE 求解：}
        欠陥補正法は $k=3$ 回の反復で CCD 演算子の $\Ord{h^6}$ 精度を完全に発現させた（§\ref{sec:verify_ppe}）．
        FD 直接法と比較して $N=128$ で $4{,}900$ 万倍高精度であり，コスト増は $3.4$ 倍に過ぎない．
        Neumann BC + ゲージ固定（§\ref{sec:verify_ppe_neumann}）では
        境界スキーム $\Ord{h^5}$ が律速する $\Ord{h^{5.0}}$ 収束を確認した．
        変密度条件（§\ref{sec:verify_ppe_varrho}）では
        滑らか密度場（$\rho_{\max}/\rho_{\min} \approx 999$）に対し
        $k = 3$ で $\Ord{h^{5.8}}$ を達成した．
        界面型密度場（smoothed Heaviside）での検証は今後の課題である．
  \item \textbf{時間積分：}
        TVD-RK3 は ODE テストで $\Ord{\Delta t^{3.00}}$ の正確な 3 次収束を確認した（§\ref{sec:verify_tvd_rk3}）．
        AB2 は同一テストで $\Ord{\Delta t^{2.0}}$ を確認し，
        初回 Euler ステップの影響は寄生根減衰により無視できることを示した（§\ref{sec:verify_ab2}）．
\end{enumerate}

\noindent\textbf{間接的に検証されるコンポーネント：}
上記のコンポーネント単体テストに加え，以下の構成要素は
本章のテストまたは第~\ref{sec:verification}章の結合テストで間接的に検証される．
\begin{itemize}[noitemsep]
  \item \textbf{Crank--Nicolson 粘性項（$\Ord{\Delta t^2}$）・ADI 分裂誤差（$\Ord{\Delta t^2}$）：}
        標準的な手法であり単体テストを省略するが，
        第~\ref{sec:verification}章の Taylor--Green 渦テスト（§\ref{sec:ns_conservation}，§\ref{sec:ns_time_accuracy}）で
        NS ソルバ全体の時間精度 $\Ord{\Delta t^2}$ として結合的に検証される．
  \item \textbf{Heaviside 正則化 $H_\varepsilon$（物性値補間）：}
        §\ref{sec:verify_field_extension}(b) の Young--Laplace テストにおいて密度場 $\rho(\psi)$ の補間が
        Laplace 圧力跳びの精度に直接反映されており，間接的に検証されている．
  \item \textbf{Block Thomas 直接解法：}
        全 CCD テスト（§\ref{sec:verify_ccd_convergence}--§\ref{sec:verify_curvature}）が
        block Thomas アルゴリズムを介して CCD 連立方程式を解いており，
        $\Ord{h^6}$ 精度の達成が block Thomas の正確性を含意する．
\end{itemize}

\begin{table}[htbp]
\centering
\caption{計算コンポーネント検証結果の総括}
\label{tab:verification_summary}
\renewcommand{\arraystretch}{1.3}
\small
\begin{tabular}{lp{0.14\linewidth}p{0.18\linewidth}p{0.24\linewidth}}
\toprule
コンポーネント & 期待精度 & 実測精度 & 備考 \\
\midrule
CCD 微分（周期） & $\Ord{h^6}$ & $\Ord{h^{6.0}}$ & $N = 256$ で丸め誤差到達 \\
CCD 微分（非一様格子） & $\Ord{h^6}$ & $\Ord{h^6}$（5.86--6.65） & GCL 合格（$5.4 \times 10^{-14}$） \\
DCCD フィルタ   & $H(\pi) = 0$ & $H(\pi) = 0$ & チェッカーボード完全抑制 \\
DCCD 移流（不連続） & TV 低減 & TV: $10.83 \to 3.15$ & CCD 比 $1/3$ \\
曲率（正弦波，$\phi$ 直接） & $\Ord{h^5}$ & $\Ord{h^{5.0}}$ & CCD 微分精度のベースライン \\
CLS 質量保存    & 保存的     & $\Ord{10^{-5}}$ ($N=256$) & 形状誤差は $\varepsilon$ が上界 \\
CLS 保存型再マッピング & LS 比改善 & 85 倍（$K=10$） & 動的非一様格子 \\
場延長・ベースライン（1D）& $\Ord{h^1}$ & $\Ord{h^{1.0}}$ & 上風差分；Hermite 補間で $\Ord{h^6}$ へ改善 \\
Hermite 延長 + CSF（Laplace） & $\Ord{h^6}$† & 相対誤差 $0.2\%$ & $N{=}128{,}256$ で単調収束；DCCD 曲率フィルタ有効 \\
PPE（欠陥補正 $k{=}3$，Dirichlet） & $\Ord{h^6}$ & $\Ord{h^{7.0}}$ & ディリクレ超収束 \\
PPE（Neumann + ゲージ固定） & $\Ord{h^5}$ & $\Ord{h^{5.0}}$ & 境界スキーム $\Ord{h^5}$ が律速 \\
PPE（変密度 $\rho_{\max}/\rho_{\min}{\approx}999$） & $\Ord{h^6}$ & $\Ord{h^{5.8}}$ & 条件数悪化でやや低下 \\
TVD-RK3         & $\Ord{\Delta t^3}$ & $\Ord{\Delta t^{3.00}}$ & ODE テストで確認 \\
AB2（初回 Euler 付き） & $\Ord{\Delta t^2}$ & $\Ord{\Delta t^{2.0}}$ & 初回 Euler の影響は寄生根で減衰 \\
\bottomrule
\multicolumn{4}{l}{\footnotesize †最近接点 Hermite 補間（§~\ref{sec:extension_pde_disc}）採用時の設計精度；
  ベースライン上風差分は $\Ord{h^1}$（表第1行参照）．} \\
\end{tabular}
\end{table}

\noindent
以上の単体テストにより，各コンポーネントが設計精度を有することを確認した．
しかし，個々の部品が正しくても，それらを\textbf{組み合わせた NS ソルバ全体}が
物理法則を正しく記述するとは限らない．
次章では，Predictor--PPE--Corrector パイプライン全体に対して
力の釣合・保存則・収束次数・安定性を検証する．
